%%%%%%%%%%%%%%%%%%%%%%%%%%%%%%%%%%%%%%%%%%%%%%%%%%%%%%%%%%%%%%%%%%%%%%%%%%%%%%%%
%%  content/04-architecture.tex — فصل ۴: معماری پیشنهادی
%%
%%  نکته: متن فارسی داخل گره‌های TikZ همیشه در \PTrl{...} پیچیده می‌شود؛ توضیحش
%%  در بخش ۵ فایل theme/persian-thesis-boxes.sty آمده است.
%%%%%%%%%%%%%%%%%%%%%%%%%%%%%%%%%%%%%%%%%%%%%%%%%%%%%%%%%%%%%%%%%%%%%%%%%%%%%%%%
\section{معماری پیشنهادی}

%%%%%%%%%%%%%%%%%%%%%%%%%%%%%%%%%%%%%%%%%%%%%%%%%%%%%%%%%%%%%%%%%%%%%%%%%%%%%%%%
\begin{frame}{معماری کلی سامانه}{جریان پردازش از راست به چپ}
  \begin{center}
    \begin{tikzpicture}[x=1cm,y=1cm]
      % ---- پنج مرحلهٔ زنجیره -------------------------------------------------
      \node[PTnodemain]   (s1) at (11.40,0) {\PTrl{پرسش کاربر}};
      \node[PTnode]       (s2) at ( 8.55,0) {\PTrl{بازیابی زندهٔ وب}};
      \node[PTnodeaccent] (s3) at ( 5.70,0) {\PTrl{بازرتبه‌بندی دومرحله‌ای}};
      \node[PTnode]       (s4) at ( 2.85,0) {\PTrl{تولید پاسخ}};
      \node[PTnodemain]   (s5) at ( 0.00,0) {\PTrl{پاسخ مستند}};

      % ---- پیکان‌ها ----------------------------------------------------------
      \draw[PTflow] (s1.west) -- (s2.east);
      \draw[PTflow] (s2.west) -- (s3.east);
      \draw[PTflow] (s3.west) -- (s4.east);
      \draw[PTflow] (s4.west) -- (s5.east);

      % ---- فناوری هر مرحله --------------------------------------------------
      \node[PTcapt,anchor=north] at ( 8.55,-0.74) {\en{Tavily Search API}};
      % «\\» بیرون از \en می‌ماند: شکستن خط کار خودِ گره است، نه کار یک بازهٔ چپ‌به‌راست.
      \node[PTcapt,anchor=north] at ( 5.70,-0.74) {\en{Cohere Rerank}\\\en{+ Jina Reranker}};
      \node[PTcapt,anchor=north] at ( 2.85,-0.74) {\en{Gemini 2.5 Flash}};

      % ---- برجسته‌سازی مشارکت اصلی ------------------------------------------
      \draw[PTaccent,line width=.7pt]
        (4.45,0.72) -- (4.45,0.96) -- (6.95,0.96) -- (6.95,0.72);
      \node[text=PTaccentdark,font=\tiny,anchor=south] at (5.70,0.99)
        {\PTrl{مشارکت اصلی این پژوهش}};
    \end{tikzpicture}
  \end{center}

  \begin{itemize}
    \item هر مرحله یک \term{واسط روشن} دارد، پس می‌توان جایگزینش کرد و اثرش را
          جداگانه سنجید.
    \item بازرتبه‌بندی، بافتار ورودی مدل را از ده‌ها سند به \hl{سه سند} کوتاه
          می‌کند.
  \end{itemize}
\end{frame}

%%%%%%%%%%%%%%%%%%%%%%%%%%%%%%%%%%%%%%%%%%%%%%%%%%%%%%%%%%%%%%%%%%%%%%%%%%%%%%%%
\begin{frame}{شرح مرحله‌به‌مرحله}
  \begin{description}
    \item[بازنویسی پرسش] پرسش کاربر به چند پرس‌وجوی جست‌وجو گسترده می‌شود تا
          یادآوری بالا برود.
    \item[بازیابی] جست‌وجوی زندهٔ وب و برداشت متن صفحات، همراه با نگه‌داشتن
          \term{نشانی منبع} برای هر قطعه.
    \item[بازرتبه‌بندی ۱] مدل \en{Cohere Rerank} فهرست خام را به ده سند برتر
          کوچک می‌کند.
    \item[بازرتبه‌بندی ۲] مدل \en{Jina Reranker} همان ده سند را به \hl{سه سند}
          نهایی می‌رساند.
    \item[تولید] مدل زبانی تنها با این سه سند پاسخ می‌سازد و در متن به آن‌ها
          ارجاع می‌دهد.
  \end{description}

  \begin{PTnote}
    دو مرحله‌ای کردن بازرتبه‌بندی، به‌جای یک مرحلهٔ بزرگ، دقت را حفظ می‌کند و
    هزینهٔ فراخوانی را پایین می‌آورد.
  \end{PTnote}
\end{frame}

%%%%%%%%%%%%%%%%%%%%%%%%%%%%%%%%%%%%%%%%%%%%%%%%%%%%%%%%%%%%%%%%%%%%%%%%%%%%%%%%
\begin{frame}{پشتهٔ پیاده‌سازی}
  \begin{center}
    \small
    \begin{tabular}{@{}ll@{}}
      \toprule
      لایه & فناوری \\
      \midrule
      کارساز و واسط برنامه‌نویسی & \en{Python 3.12} + \en{FastAPI} \\
      پایگاه داده                & \en{SQLite} + \en{SQLAlchemy} \\
      واسط کاربری                & \en{HTML5} + \en{Tailwind CSS v4} \\
      مدل زبانی                  & \en{Gemini 2.5 Flash} \\
      جست‌وجوی وب                & \en{Tavily Search API} \\
      بازرتبه‌بندی               & \en{Cohere Rerank}، \en{Jina Reranker} \\
      \bottomrule
    \end{tabular}
  \end{center}

  \begin{itemize}
    \item همهٔ سرویس‌های بیرونی پشت یک \term{لایهٔ واسط} قرار دارند؛ تعویض هر یک
          فقط یک کلاس را تغییر می‌دهد.
  \end{itemize}
\end{frame}

%%%%%%%%%%%%%%%%%%%%%%%%%%%%%%%%%%%%%%%%%%%%%%%%%%%%%%%%%%%%%%%%%%%%%%%%%%%%%%%%
%%  اسلاید کد
%%  ۱) هر اسلایدی که PTcode دارد باید [fragile] بگیرد.
%%  ۲) داخل PTcode فقط نویسهٔ اَسکی بنویسید. قلم تک‌عرض پیش‌فرض نشانهٔ فارسی ندارد
%%     و listings هم با نویسه‌های چندبایتی خوب کنار نمی‌آید؛ پس توضیحات کد را
%%     انگلیسی بنویسید و شرح فارسی را در فهرست زیر اسلاید بیاورید.
%%%%%%%%%%%%%%%%%%%%%%%%%%%%%%%%%%%%%%%%%%%%%%%%%%%%%%%%%%%%%%%%%%%%%%%%%%%%%%%%
\begin{frame}[fragile]{هستهٔ زنجیره}{ساده‌سازی‌شده برای ارائه}
\begin{PTcode}[language=Python]
async def answer(question: str) -> Answer:
    queries   = rewrite(question)              # query expansion
    documents = await search_web(queries)      # Tavily
    top10     = cohere_rerank(question, documents, k=10)
    top3      = jina_rerank(question, top10,   k=3)
    return generate(question, context=top3)    # Gemini 2.5 Flash
\end{PTcode}

  \begin{itemize}
    \item کل زنجیره در یک تابع ناهم‌زمان جمع می‌شود؛ هر گام یک وابستگی
          جایگزین‌پذیر است.
    \item دو فراخوانی \en{rerank} پشت سر هم می‌آیند: نخست ده سند، سپس سه سند.
  \end{itemize}
\end{frame}
